\chapter{Architecture}
\label{ch:arch}

\section{Roles}

\begin{table}[H]
\centering
\begin{tabularx}{\textwidth}{L{0.24\textwidth} Y}
\toprule
\textbf{Role} & \textbf{Daemons} \\
\midrule
Entry node & \kn{master}, \kn{collector}, \kn{negotiator}, \kn{schedd}, \kn{startd} (limited), \kn{shared\_port} \\
Worker & \kn{master}, \kn{startd} \\
User PC (desktop) & Worker daemons plus the submit client (\kn{condor\_submit -remote -spool}); \textbf{no schedd} \\
Submit-only client & \textbf{No daemons at all}; submit client only. Typically laptops \\
\bottomrule
\end{tabularx}
\caption{Daemon layout by role.}
\end{table}

\section{Not every user machine is a worker}
\label{sec:notallworkers}

Donation is expected of \textbf{desktops}, which sit still and stay plugged in.
A laptop that only submits is a normal, supported member of the pool.

\begin{table}[H]
\centering
\begin{tabularx}{\textwidth}{Y L{0.18\textwidth} L{0.24\textwidth}}
\toprule
& \textbf{Desktop} & \textbf{Laptop} \\
\midrule
Reachable inbound by the entry node & yes & only while on the LAN \\
Awake while a job runs & yes & the lid decides \\
Owner's battery and fans & mains & theirs \\
\bottomrule
\end{tabularx}
\caption{Why the execute role does not travel.}
\end{table}

A worker must be \textbf{reachable inbound} to be claimed, so an off-LAN laptop
would register, match, and then fail the claim. A submit client makes only
\textbf{outbound} connections, which is why it works through a tunnel, behind a
firewall, and from outside the office.

Take execute capacity from machines that sit still. Take submissions from
anything.

There is exactly one schedd, on the entry node. This is deliberate: a schedd on
a laptop dies with the lid, taking its queue with it. Centralising it means a
user's jobs survive their machine going to sleep, which is the normal state of a
laptop between two meetings.

\begin{keypoint}
A user PC runs worker daemons \emph{and} the submit client, but no schedd.
Submission is a remote operation against the entry node's queue; the local
Condor installation exists only to donate the machine's capacity.
\end{keypoint}

\section{Choosing the entry node}

The selected class is a \textbf{low-power four-core mini-PC}, not the pool's
highest-core machine. That is counter-intuitive enough to justify.

\begin{table}[H]
\centering
\begin{tabularx}{\textwidth}{L{0.22\textwidth} Y}
\toprule
\textbf{Criterion} & \textbf{Finding} \\
\midrule
Role bottleneck & \kn{fsync} latency on \file{job\_queue.log}, RAM per shadow, and NIC bandwidth --- not CPU \\
Core count & Sixteen threads buy nothing for this role, and the high-core machine is the pool's only one \\
Commodity & The four-core class is the fleet's repeated unit; the flagship is scarce \\
Test validity & The production entry node will be commodity class. Testing on the flagship would teach nothing transferable \\
Scaling path & Horizontal --- add schedds, or split collector and negotiator off. Never vertical \\
\bottomrule
\end{tabularx}
\caption{Entry node selection criteria.}
\end{table}

Spending the pool's only sixteen-thread machine on a role bounded by disk
latency would waste its one genuinely scarce resource. It is worth far more as a
worker, and that is where it sits.

Requirements for whichever machine holds the role: \kn{SPOOL} on NVMe, and about
half its cores donated to compute with the remainder reserved for the daemons.

\section{Machine classes}

\begin{table}[H]
\centering
\begin{tabularx}{\textwidth}{L{0.27\textwidth} l l l l l}
\toprule
\textbf{Class} & \textbf{Arch} & \textbf{Cores} & \textbf{RAM} & \textbf{Disk free} & \textbf{Count} \\
\midrule
Mini-PC (entry + worker) & x86\_64  & 4   & 16\,G & 30--70\,G & 8+, growing \\
Flagship worker          & x86\_64  & 16  & 31\,G & $\sim$270\,G & 1 \\
Laptop                   & x86\_64  & 12  & 8\,G  & $\sim$250\,G & 1+ \\
Edge SBC                 & aarch64  & 4   & 4\,G  & $\sim$36\,G (SD) & 1+ \\
Jetson                   & aarch64  & TBD & TBD   & TBD & future, GPU \\
\bottomrule
\end{tabularx}
\caption{Hardware classes in the fleet.}
\end{table}

Distributions: Ubuntu 24.04 LTS on x86\_64, Debian 13 on the edge SBC, and L4T
expected on the Jetsons. Three distributions and two architectures is the point
at which ``just ship the binary'' stops working; Chapter~\ref{ch:roadmap}
returns to it.

\section{The submission path}

\begin{lstlisting}[style=term]
user's own PC                    entry node                    worker
-------------                    ----------                    ------
condor_submit -remote -spool ->  schedd accepts job
                                 input files land in SPOOL
                                 negotiator matches       ->    startd claims
                                 shadow forks             <->   starter runs job
                                 output returns to SPOOL
condor_transfer_data          <- output released
\end{lstlisting}

Two properties of this path matter more than they first appear.

\textbf{The user needs no account anywhere except an authentication token.}
Submission is authenticated with IDTOKENS over a single TCP port; there is no
shell, no home directory and no shared filesystem involved.

\textbf{All job I/O funnels through the shadows on the entry node.} That is the
pool's scaling limit --- not CPU, not RAM. Chapter~\ref{ch:donation} sizes the
disk defences around it, and the answer for genuinely large data is shared
storage with a matching \kn{FILESYSTEM\_DOMAIN} so that matched jobs skip
transfer entirely, not a bigger entry node.

\section{Workload constraints}

Eviction is the normal case in this design, not an exception, and users have to
be told so up front.

\begin{itemize}
  \item Many short jobs, not few long ones. A ten-hour job on a desktop may never complete.
  \item Jobs must be idempotent --- restartable from scratch without corrupting state.
  \item Long work uses self-checkpointing (\kn{checkpoint\_exit\_code}).
  \item ``My job keeps restarting'' is a design property of a desktop pool, not a bug.
\end{itemize}

\section{Scheduling and fairness}

\begin{table}[H]
\centering
\begin{tabularx}{\textwidth}{L{0.30\textwidth} Y}
\toprule
\textbf{Mechanism} & \textbf{Purpose} \\
\midrule
\kn{condor\_userprio} fair-share & Heavy recent users decay in priority automatically \\
Machine \kn{RANK} & Owner preference on their own hardware \\
\kn{NEGOTIATOR\_PRE\_JOB\_RANK} & Fill the entry node last \\
\kn{MAX\_JOBS\_RUNNING} & Ceiling on concurrent shadows; converts ``entry node falls over'' into ``jobs stay queued'' \\
\kn{MAX\_TRANSFER\_INPUT\_MB} & Refuse fat submissions rather than filling spool \\
\bottomrule
\end{tabularx}
\caption{Fairness and self-defence mechanisms.}
\end{table}

\kn{RANK = (RemoteOwner =?= "<owner>")} gives each owner first call on their own
machine, including preempting a stranger's job. This is the adoption lever. The
failure mode it defends against is quiet: nobody files a bug when Condor makes
their PC stutter --- they run \kn{systemctl disable condor}, and the pool
shrinks without anyone noticing.
